% !PRINTED 348
\weberpage{348}
%<<CONT
arises, $(1,1)$, $(2,2)$ can still pass over into $(0,2)$, $(0,1)$, but not into
$(2,0)$, $(1,0)$ or into $(1,2)$, $(2,1)$. Hence the group $S_3 Q$ is likewise not yet
doubly transitive, but it is primitive. If, however, in order to form $S_2$,
we adjoin the substitution

\[
\begin{pmatrix} 0, & 1 \\ -1, & 0 \end{pmatrix}
\]

then $(1,1)$, $(2,2)$ passes over into $(0,2)$, $(0,1)$ and into $(1,2)$, $(2,1)$ and
$(2,0)$, $(1,0)$. Hence the group $S_2 Q$ is doubly transitive, and consequently so are
$S_1 Q$ and $S Q$. The triple equations therefore first begin with $S_2 Q$.
The equations with narrower groups are not yet proper triple equations.

We can form still other linear groups contained in $S Q$, e.\,g. if we
denote by $S'$ the cyclic group of degree 3:

\[
S' = \begin{pmatrix} 1, & 0 \\ 0, & 1 \end{pmatrix}, \begin{pmatrix} 1, & 0 \\ 1, & 1 \end{pmatrix}, \begin{pmatrix} 1, & 0 \\ -1, & 1 \end{pmatrix}
\]

the group of degree 27, $S' Q = Q S'$. Through the group $S'$,
$(1,1)$ $(2,2)$ passes over into $(1,1)$ $(2,2)$, $(1,2)$ $(2,1)$, $(1,0)$ $(2,0)$, while
$(0,1)$ $(0,2)$ does not arise from it. Hence this too is not the group of a
triple equation.

Here there would still be a region for further investigations into these special
metacyclic equations of degree 9 and the representation of their
roots according to the principles of \gesperrt{Kronecker} (Bd. I, eighteenth chapter).

\begin{center}
\S.\,94.\\[2pt]
\gesperrt{Reality Relations of the Triple Equations.}
\end{center}

If the domain of rationality is real, we can draw some general conclusions
about the reality of the roots of a triple equation of degree 9.
For if in this case two real roots occur in one triple, then,
as follows from the equation $x_3 = \Theta (x_1, x_2)$, the third root must also be
real. If, therefore, there is any imaginary root at all, then in each of the
four triples to which this root belongs there must occur at least one further imaginary
root, and consequently the number of imaginary roots is at least equal to 5,
or, since the number of imaginary roots
%CONT>>
